\section{Cómo convertir Llama en un Transformer de generación de video}

Solo después de fijar la representación y el objetivo de aprendizaje llega el turno del backbone. Movie Gen no transfiere directamente los pesos de un modelo de lenguaje al video, sino que elige los bloques Llama-style y la infraestructura de entrenamiento que ya conoce, y entrena desde una inicialización aleatoria. Para adaptarse a Flow Matching condicional, agrega cross-attention de texto y AdaLN de timestep, y sustituye la causal GQA por MHA bidireccional completa.

\subsection{Por qué elegir la forma de Llama y no los pesos de Llama}

El criterio de ingeniería del ponente es «elegir una forma que ya se sabe cómo escalar». Los kernels, las estrategias de paralelismo, el monitoreo y la experiencia con fallas del Transformer Llama-style ya están maduros; el valor de esa infraestructura puede superar la ventaja local de algún módulo visual especializado. La figura de la llama que sigue es una metáfora de arquitectura, no evidencia de transferencia de conocimiento lingüístico.

\lecturefigure{slide-193-llama-backbone.jpg}{El latent del video se divide en patches y se envía a un Transformer Llama-style con inicialización aleatoria}{00:27:15}

\begin{warningbox}{Movie Gen no es «hacer fine-tuning de Llama para obtener un modelo de video»}
Este proyecto usa un backbone Llama-style con inicialización aleatoria. En clase alguien preguntó si se había inicializado a partir de pretrained Llama, y el docente respondió con claridad que eso no se intentó en este proyecto. Si es posible transferir entre modalidades distintas y objetivos distintos sigue siendo una pregunta de investigación abierta.
\end{warningbox}

\teachervoice{00:25:12--00:27:15, el docente explica que la razón principal para elegir la forma de Llama es que la ruta de escalamiento y la infraestructura ya son conocidas; esto no significa que se hayan usado en Movie Gen los pesos preentrenados con texto.}

\subsection{Modificación uno: el texto entra mediante cross-attention}

La sección anterior eligió la estructura principal Llama-style; en esta se empiezan a reincorporar, una por una, las condiciones de generación. Los tokens de video y los de texto tienen longitudes distintas, y al texto no hace falta añadirle ruido como en Flow Matching; el problema central es cómo hacer que cada posición del video lea el significado de las palabras relevantes sin comprimir un prompt largo en un solo vector. Movie Gen usa tres encoders de texto complementarios para producir la representación condicional, y luego deja que el flujo de video la lea mediante cross-attention.

\lecturefigure{slide-194-text-conditioning.jpg}{Las tres representaciones de texto se concatenan y condicionan el Transformer de video mediante cross-attention}{00:29:23}

Para los hidden states del video $H$ y la representación del texto $P$, la cross-attention puede escribirse como
\[
Q=HW_Q,\quad K=PW_K,\quad V=PW_V,
\qquad
\operatorname{CA}(H,P)=\operatorname{softmax}\!\left(\frac{QK^\top}{\sqrt d}\right)V.
\]
$W_Q,W_K,W_V$ son matrices de proyección y $d$ es la head dimension. La longitud de la salida sigue siendo igual al número de tokens de video; el texto actúa como memory condicional y la misma capa no le añade ruido en sentido inverso.

\begin{knowledgebox}{Por qué se necesitan varios text encoders}
Cada encoder destaca en una escala distinta: uno aporta semántica global, otro conserva prompts largos y relaciones a nivel de palabra, y otro está alineado con la semántica visual. Concatenar varias representaciones aumenta la capacidad condicional, pero también la memoria de GPU, la latencia y la complejidad de la interfaz; su beneficio debe verificarse con estudios de ablación.
\end{knowledgebox}

\subsection{Modificación dos: AdaLN escribe el paso de tiempo en cada block}

La ruta del texto ya responde «qué generar»; esta sección pasa a «en qué punto de la eliminación de ruido se está». Estos dos tipos de condición tienen estructuras de información distintas, por lo que no deben compartir mecánicamente la misma interfaz. Movie Gen retoma de DiT la Adaptive LayerNorm (AdaLN, normalización de capa adaptativa): el embedding del tiempo produce scale y shift, de modo que el mismo block ejecuta cálculos distintos según el nivel de ruido.

\lecturefigure{slide-195-adaptive-layernorm.jpg}{AdaLN usa el timestep para generar el scale y el shift de cada capa}{00:30:04}

Si $h$ es el hidden state de un token y $e_t$ es el embedding del tiempo, entonces
\[
\operatorname{AdaLN}(h;t)
=\gamma(e_t)\odot\operatorname{LN}(h)+\beta(e_t),
\]
donde $\gamma,\beta$ son salidas de redes pequeñas y $\odot$ es la multiplicación elemento a elemento. El costo de AdaLN no crece con la longitud del texto, por lo que es adecuada para una condición de timestep de baja dimensión; el texto sigue conservando su estructura de secuencia mediante cross-attention.

\begin{importantbox}{Las interfaces de condición se reparten según la estructura de la información}
El texto en forma de secuencia larga requiere una lectura token-to-token selectiva, por eso pasa por cross-attention; una sola variable de tiempo continua solo necesita modular las estadísticas de cada capa, por eso pasa por AdaLN. Comprimir todas las condiciones en un vector o concatenarlas todas como tokens haría perder esta división estructurada del trabajo.
\end{importantbox}

\subsection{Modificación tres: de causal GQA a bidirectional MHA}

Un modelo de lenguaje solo puede ver los tokens a su izquierda porque su objetivo es predecir la siguiente palabra; Flow Matching actualiza al mismo tiempo todo el noisy latent, así que no existe un orden causal. Movie Gen usa full bidirectional attention y elige MHA en lugar de GQA, de modo que cada query head conserve sus propias key/value heads.

\lecturefigure{slide-196-bidirectional-attention.jpg}{Flow Matching de video usa atención bidireccional, no la causal mask de los modelos de lenguaje}{00:30:38}

\begin{knowledgebox}{Primera aparición: MHA y GQA}
MHA (Multi-Head Attention) asigna K/V heads independientes a cada query head; GQA (Grouped-Query Attention) hace que varias query heads compartan menos K/V heads, principalmente para reducir la KV cache autoregresiva. Movie Gen no genera de forma autoregresiva token por token y durante el entrenamiento tampoco depende de una KV cache de largo plazo, por lo que adopta MHA bidireccional, que tiene mayor capacidad expresiva.
\end{knowledgebox}

\subsection{Arquitectura completa y el costo de sistema de 73K tokens}

Solo al encadenar TAE, patchify, ruido, condición de texto, timestep y decoder se obtiene el sistema de generación completo. Al leer la figura conviene seguir la forma de los datos de izquierda a derecha: el video se codifica primero como latent, el latent y el ruido gaussiano forman $X_t$, el Transformer predice la velocidad y, por último, el decoder de TAE devuelve el video en píxeles.

\lecturefigure{slide-244-full-architecture.jpg}{Arquitectura completa de Movie Gen Video: TAE, patchify, Transformer condicional y decoder}{00:35:12}

El término dominante de la atención completa es aproximadamente
\[
\operatorname{FLOPs}_{\mathrm{attn}}=O(LN^2d),
\qquad
\operatorname{Memory}_{\mathrm{attn}}=O(N^2),
\]
donde $L$ es el número de capas, $N$ la longitud de la secuencia y $d$ la head dimension. Con $N=73{,}000$, un solo attention map completo tiene unos $5.3\times10^9$ elementos y no puede almacenarse de forma ingenua en una sola GPU.

Por ello el entrenamiento requiere varios tipos de paralelismo: data parallel replica el modelo y procesa las muestras por lotes; tensor parallel divide las multiplicaciones de matrices; sequence/context parallel fragmenta a lo largo de la dimensión de tokens. Aquí sharding (fragmentación) significa repartir parámetros, activaciones o secuencias entre varias GPU, de modo que cada una guarde solo una parte en lugar de una copia completa. Las collectives (comunicación colectiva) incluyen all-reduce, all-gather, reduce-scatter y otras primitivas de comunicación entre varias GPU, que se usan para agregar o intercambiar resultados fragmentados.

\begin{warningbox}{El número de parámetros no es la única escala del sistema}
Los 30B parámetros describen la capacidad de los pesos, el context de 73K describe las activaciones y la attention de una sola muestra, y el número de pasos de muestreo describe cuántas veces se repite la pasada hacia adelante. El costo de despliegue debe calcular a la vez la memoria de GPU de los pesos, las activaciones, los buffers temporales, la comunicación colectiva y los ODE steps; reportar solo el número de parámetros subestima gravemente la generación de video.
\end{warningbox}

\teachervoice{00:33:59--00:35:12, el docente atribuye el progreso de dos años al cambio conjunto de arquitectura, datos, escala del modelo, compute y hardware. El valor de la arquitectura completa no está en un block nuevo, sino en que todas estas partes por fin pueden escalar juntas dentro de un sistema unificado.}

\subsection{Resumen de la sección}

Movie Gen usa un backbone Llama-style con inicialización aleatoria, no una transferencia de pesos de lenguaje. Las tres modificaciones principales son la cross-attention de texto, la AdaLN de timestep y la MHA bidireccional. El verdadero cuello de botella del sistema completo proviene de la secuencia de 73K tokens, la attention cuadrática, la comunicación de la fragmentación entre varias GPU y la inferencia ODE de múltiples pasos, no solo de los 30B parámetros.
